\documentclass[11pt]{article}
\usepackage{reusv}
\usepackage{booktabs}
\usepackage{amsmath,amssymb}
\usepackage{siunitx}
\usepackage{graphicx}
\usepackage{multirow}
\usepackage{xcolor}
\usepackage{listings}

% ── Python 코드 스타일 ──
\definecolor{codebg}{HTML}{F5F5F0}
\definecolor{codeframe}{HTML}{CCCCBB}
\definecolor{codekw}{HTML}{0060A0}
\definecolor{codecomment}{HTML}{60802A}
\definecolor{codestring}{HTML}{B03030}
\lstdefinestyle{pythonstyle}{
  language=Python,
  backgroundcolor=\color{codebg},
  frame=single,
  rulecolor=\color{codeframe},
  framesep=6pt,
  xleftmargin=8pt,
  xrightmargin=8pt,
  basicstyle=\small\ttfamily,
  keywordstyle=\color{codekw}\bfseries,
  commentstyle=\color{codecomment}\itshape,
  stringstyle=\color{codestring},
  showstringspaces=false,
  breaklines=true,
  tabsize=4,
  numbers=left,
  numberstyle=\tiny\color{gray},
  numbersep=6pt,
  aboveskip=10pt,
  belowskip=6pt,
}

\setborderthickness{0.5pt}
\setbordermargin{1cm}
\setbordercolor{black}

\slimtitle{드리프트 적분의 솔버 통합과 성능 최적화}

\begin{document}

\drawborder
\thispagestyle{firstpage}

\lightertitle{드리프트 적분의 솔버 통합과 성능 최적화}

{\normalsize\textit{Research Note No.~011 $|$ bezier-orbit-transfer-scp}}\par\medskip

\def\lighterdate{March 28, 2026}
\lighterauthor{Suwon Lee$^{\dagger}$}

\lighteraddress{$^\dagger$}{Future Mobility Control Lab, Kookmin University, Seoul, Korea}
\lighteremail{suwon.lee@kookmin.ac.kr}

%% ── 초록 ─────────────────────────────────────────────────────
\begin{abstract}
보고서 009(아핀 근사)\cite{report009}와 010(Bernstein 대수)\cite{report010}에서
제안한 드리프트 적분 가속 방법론을 SCP 솔버에 통합하고, 구현 과정에서 발견된
자기일관 중력 보정 문제를 해결한다.
\texttt{bern\_product}/\texttt{bern\_compose}의 벡터화, 가중치 행렬 캐싱,
워밍스타트, 적응적 조기 종료, 보정용 $K$ 분리를 적용하여
초기 Python 구현 대비 \textbf{180배} 속도 향상을 달성하였다.
최종 벤치마크에서 아핀 방법(수치 보정)은 RK4 대비 \textbf{1.97배},
아핀 방법(대수 보정)은 \textbf{1.33배} 빠르면서 유사한 비용과 수렴을 보인다.
마지막으로, 대수적 보정이 연산 속도보다는 \textbf{제어점 최적화 구조의 보존}에
본질적 가치를 가짐을 논의한다.
\end{abstract}

%% ══════════════════════════════════════════════════════════════
\section{도입}
%% ══════════════════════════════════════════════════════════════

SCP 기반 궤도전이 최적화에서, 매 반복마다 비선형 중력의 드리프트 적분
\begin{align}
  \mathbf{c}_v &= \int_0^1 \mathbf{a}_{\mathrm{grav}}\!\bigl(\mathbf{r}_{\mathrm{ref}}(\tau)\bigr)\,d\tau, \label{eq:cv}\\
  \mathbf{c}_r &= \int_0^1\!\int_0^s \mathbf{a}_{\mathrm{grav}}\!\bigl(\mathbf{r}_{\mathrm{ref}}(\sigma)\bigr)\,d\sigma\,ds \label{eq:cr}
\end{align}
을 계산해야 한다\cite{report006}.
기존 구현은 RK4 수치 전파로 참조 궤적을 얻은 뒤 사다리꼴 적분을 사용하며,
보고서~008\cite{report008}의 프로파일링에 따르면 이 과정이 SCP 반복 비용의
86\%를 차지한다.

보고서~009\cite{report009}는 구간별 Taylor 선형화(아핀 근사)로,
보고서~010\cite{report010}은 Bernstein 대수 파이프라인으로
RK4를 대체하는 이론을 제시하였으나, 두 방법 모두 솔버에 통합되지 않은
상태였다.
본 보고서에서는 (1) 세 방법을 SCP 솔버에 통합하고,
(2) 구현 과정에서 발견된 자기일관 중력 보정 문제를 해결하며,
(3) Python 레벨 최적화를 통해 대수적 방법의 실용적 속도를 달성하고,
(4) 대수적 방법이 갖는 최적화 구조적 함의를 논의한다.


%% ══════════════════════════════════════════════════════════════
\section{솔버 통합 설계}
%% ══════════════════════════════════════════════════════════════

\subsection{2축 분류: 드리프트 방법 $\times$ 중력 보정 모드}

드리프트 적분 계산을 두 가지 직교 축으로 분류한다.

\paragraph{축 1: 드리프트 적분 방법 (\texttt{method})}
보정된 참조 위치에서 $\mathbf{c}_v$, $\mathbf{c}_r$을
\emph{어떻게 계산하느냐}의 차이이다.

\begin{itemize}
  \item \textbf{RK4} — 기존 방식. RK4 전파 + 사다리꼴 적분.
  \item \textbf{아핀 근사} — $[0,1]$을 $K$개 구간으로 분할,
        구간 중점에서 중력 평가, midpoint rule\cite{report009}.
  \item \textbf{Bernstein 대수} — \texttt{gravity\_composition\_pipeline}으로
        중력 가속도 제어점을 구한 뒤, 정적분 성질
        $\int_0^1 p(\tau)\,d\tau = \text{mean}(\mathbf{p})$를 사용\cite{report010}.
\end{itemize}

\paragraph{축 2: 중력 보정 모드 (\texttt{gravity\_correction})}
참조 위치 곡선 자체를 \emph{어떻게 구성하느냐}의 차이이다.
둘 다 ``제어 전용 위치에 중력 이중 적분 $\Phi(\tau)$를 더한다''는
목적은 같다 (\S\ref{sec:selfconsistent} 참조).

\begin{itemize}
  \item \textbf{대수적 (\texttt{algebraic})} — Bernstein 파이프라인으로
        중력 가속도 제어점을 구한 뒤, 적분 행렬 곱으로 이중 적분.
        수치적분 불사용.
  \item \textbf{수치적 (\texttt{numerical})} — 이산 점에서 중력 평가,
        사다리꼴 이중 적분, 최소자승 Bernstein 적합.
\end{itemize}

두 축이 독립적이므로 $3 \times 2 = 6$가지 조합이 가능하다
(RK4는 자체 전파를 사용하므로 보정 모드 불필요, 실질 5가지).

\subsection{\texttt{DriftConfig} 데이터클래스}

선택은 \texttt{DriftConfig}로 캡슐화된다:

\begin{lstlisting}[style=pythonstyle]
DriftConfig(
    method="bernstein",       # "rk4"|"affine"|"bernstein"
    gravity_correction="algebraic", # "algebraic"|"numerical"
    bernstein_K=8,            # pipeline approx degree
    bernstein_R=20,           # degree reduction target
    affine_K=8,               # affine segment count
    gravity_correction_K=4,   # correction pipeline K
    n_gravity_iter=8,         # max self-consistent iter
    gravity_tol=1e-6,         # early-stop tolerance
)
\end{lstlisting}

\subsection{Inner loop 분기 로직}

SCP 반복 루프 내에서 \texttt{compute\_drift} 디스패처가 방법에 따라 분기한다.
비RK4 방법에서는 반복 중 RK4 전파를 완전히 건너뛰어 비용을 절감하되,
\textbf{수렴 후 최종 검증은 항상 고해상도 RK4}로 수행하여 BC 위반을
정확히 측정한다.


%% ══════════════════════════════════════════════════════════════
\section{자기일관 중력 보정}
\label{sec:selfconsistent}
%% ══════════════════════════════════════════════════════════════

\subsection{문제: 제어 전용 위치의 한계}

아핀·Bernstein 방법에서 참조 위치는 적분 체인으로 구성된다:
\begin{equation}
  \mathbf{q}_r = \mathbf{r}_0 + t_f \mathbf{v}_0 \boldsymbol{\ell}
               + t_f^2 \bar{I}_N \mathbf{P}_u
  \label{eq:qr_control}
\end{equation}
여기서 $\boldsymbol{\ell} = [0, \frac{1}{N+2}, \ldots, 1]^T$이다.
이 위치에는 \textbf{중력의 위치 기여}가 포함되어 있지 않다.

LEO 궤도에서 $t_f \approx 4\,\text{TU}$ 비행 시, 위성은 중력에 의해
궤도를 거의 한 바퀴 선회한다.
식~\eqref{eq:qr_control}의 위치는 이를 반영하지 못하여
실제 궤적과 크게 괴리되며, 이 위치에서 중력을 평가하면
SCP가 발산한다.

수치 실험에서, 중력 보정 없이 Bernstein/아핀 방법을 적용하면
$\tau = 0.5$에서의 위치 오차가 $\sim1.7$ (정규화 단위)에 달하고,
BC 위반이 $\sim5.5$로 수렴하지 못함을 확인하였다.

\subsection{해법: 자기일관 반복}

중력의 이중 적분
\begin{equation}
  \Phi(\tau) = \int_0^\tau\!\int_0^s
    \mathbf{a}_{\mathrm{grav}}\!\bigl(\mathbf{r}(\sigma)\bigr)\,d\sigma\,ds
  \label{eq:phi}
\end{equation}
를 반복적으로 위치에 더하여 실제 궤적에 가까운 참조를 구성한다:
\begin{equation}
  \mathbf{q}_r^{(j+1)} = \mathbf{q}_{r,\mathrm{control}}
    + t_f^2 \, \Phi^{(j)}
  \label{eq:selfconsistent}
\end{equation}

\paragraph{대수적 보정.}
\texttt{gravity\_composition\_pipeline}으로 중력 가속도 제어점
$\mathbf{q}_a \in B_R$을 구한 뒤, 적분 행렬 곱으로 이중 적분 제어점을 얻는다:
\begin{equation}
  \mathbf{q}_\Phi = I_{R+1} \cdot I_R \cdot \mathbf{q}_a
  \quad \in B_{R+2}
  \label{eq:phi_algebraic}
\end{equation}
이후 $L^2$ 최적 차수 축소로 $B_{N+2}$에 사영하여 위치에 더한다.
모든 연산이 제어점 공간의 행렬 곱으로 수행된다.

\paragraph{수치적 보정.}
$n_{\mathrm{pts}}$개 점에서 중력을 평가하고, 사다리꼴 이중 적분으로
$\Phi(\tau)$를 구한 뒤, 최소자승으로 Bernstein 제어점에 적합한다.

\paragraph{보정 결과.}
8회 반복 후 $\tau = 0.5$에서의 위치 오차가 $1.69 \to 0.015$
(100배 이상 개선)으로 줄어, SCP가 정상 수렴하게 된다.


%% ══════════════════════════════════════════════════════════════
\section{벡터화 최적화}
\label{sec:vectorization}
%% ══════════════════════════════════════════════════════════════

초기 구현에서 대수적 보정은 RK4 대비 100배 이상 느렸다.
원인 분석과 최적화 과정을 기술한다.

\subsection{병목 진단: Python 루프 오버헤드}

\texttt{bern\_product}의 초기 구현은 이중 Python 루프였다:
\begin{lstlisting}[style=pythonstyle]
for k in range(N + M + 1):
    for j in range(j_min, j_max + 1):
        r[k] += C(N,j)*C(M,k-j)/C(N+M,k) * p[j] * q[k-j]
\end{lstlisting}
$N = M = 28$일 때 $\sim$588회의 Python 루프 반복이 발생하며,
반복당 함수 호출·인덱싱 오버헤드가 $\sim$1--5\,$\mu$s이므로
총 $\sim$3\,ms가 소요된다.
\texttt{bern\_compose} ($K=8$)는 이를 $\sim$24회 호출하므로 $\sim$100\,ms이다.

실제 부동소수점 연산(FLOP) 수는 $\sim$25,000으로, numpy 벡터 연산으로는
$\sim$0.003\,ms에 처리 가능하다.
즉 \textbf{40,000배}의 오버헤드가 Python 루프에서 발생하고 있었다.

\subsection{적용된 최적화}

\begin{enumerate}
  \item \textbf{가중치 행렬 캐싱}:
    $W_{N,M}[k,j] = C(N,j) C(M,k\!-\!j) / C(N\!+\!M,k)$를
    \texttt{@lru\_cache}로 $(N,M)$ 쌍별로 사전 계산.
    Toeplitz 인덱스 구조도 동일하게 캐싱.

  \item \textbf{\texttt{bern\_product} 벡터화}:
    이중 루프를 단일 numpy 연산으로 대체:
\begin{lstlisting}[style=pythonstyle]
r = np.sum(W * p[np.newaxis, :] * Q_shifted, axis=1)
\end{lstlisting}

  \item \textbf{\texttt{degree\_elevate} 벡터화}:
    내부 루프를 배열 슬라이스 연산으로 대체.

  \item \textbf{\texttt{bern\_compose} 개선}:
    $g^i \cdot (1\!-\!g)^{K-i}$의 차수가 항상 $KM$이므로
    불필요한 \texttt{degree\_elevate} 호출 제거.

  \item \textbf{\texttt{degree\_reduce} 벡터화}:
    교차 Gram 행렬 $C_{Q \to R}$을 \texttt{@lru\_cache}로 캐싱.
    2D 배열 입력 시 한 번의 \texttt{np.linalg.solve}로 처리.

  \item \textbf{Chebyshev 근사 캐싱}:
    $h(s) = s^{-3/2}$의 $K$차 Bernstein 근사를
    반올림된 $[a,b]$ 구간 키로 캐싱.

  \item \textbf{보정용 $K$ 분리}:
    자기일관 보정에는 \texttt{gravity\_correction\_K}$=4$를 사용하여
    \texttt{bern\_compose} 비용을 절반으로 줄임.
    최종 드리프트 계산은 \texttt{bernstein\_K}$=8$로 정확도 유지.

  \item \textbf{워밍스타트 + 적응적 조기 종료}:
    이전 SCP 반복의 $\Phi$ 제어점을 재사용하여 초기값으로 제공.
    $\|\Phi^{(j+1)} - \Phi^{(j)}\| < \texttt{gravity\_tol}$ 시 조기 종료.
    SCP 수렴 근방에서 보정 반복이 8회 $\to$ 1--2회로 감소.
\end{enumerate}

\subsection{최적화 단계별 성능 변화}

표~\ref{tab:optim_stages}에 시나리오~B (LEO$\to$1.2$a_0$)에서의
SCP 전체 풀이 시간 변화를 정리한다.

\begin{table}[h]
\centering
\caption{최적화 단계별 SCP 전체 풀이 시간 (Affine, 대수적 보정)}
\label{tab:optim_stages}
\begin{tabular}{lrrr}
\toprule
단계 & 시간 [ms] & RK4 대비 & 누적 배속 \\
\midrule
초기 (Python 루프)    & 25,564 & 0.01$\times$ & 1$\times$ \\
벡터화 적용 후        & 392    & 0.48$\times$ & 65$\times$ \\
+ 워밍스타트/조기종료  &        &              & \\
\quad + $K_\mathrm{corr}=4$ &   &              & \\
\quad + Chebyshev 캐싱 & 142   & 1.33$\times$ & 180$\times$ \\
\midrule
RK4 (비교)            & 189   & 1.00$\times$ & --- \\
\bottomrule
\end{tabular}
\end{table}


%% ══════════════════════════════════════════════════════════════
\section{성능 벤치마크}
%% ══════════════════════════════════════════════════════════════

\subsection{SCP 전체 풀이}

표~\ref{tab:scp_final}은 최종 최적화 후의 SCP 전체 풀이 결과이다.
시나리오~B (LEO$\to$1.2$a_0$, $N=12$, $t_f^* = 3.62$),
$\texttt{tol\_ctrl} = 10^{-3}$, $\texttt{tol\_bc} = 0.05$.

\begin{table}[h]
\centering
\caption{SCP 전체 풀이 최종 벤치마크 (시나리오 B)}
\label{tab:scp_final}
\begin{tabular}{lrrrcrr}
\toprule
방법 & 시간 & 배속 & 비용 $J$ & 수렴 & 반복 & BC 위반 \\
\midrule
RK4              & 189\,ms  & 1.00$\times$ & 0.00665 & $\checkmark$ & 24 & $3.3 \!\times\! 10^{-3}$ \\
Affine(num)      & 96\,ms   & 1.97$\times$ & 0.00653 & $\checkmark$ & 23 & $2.4 \!\times\! 10^{-2}$ \\
Affine(alg)      & 142\,ms  & 1.33$\times$ & 0.00621 & $\checkmark$ & 23 & $3.1 \!\times\! 10^{-2}$ \\
Bern(num)        & 150\,ms  & 1.26$\times$ & 0.00635 & $\checkmark$ & 29 & $3.3 \!\times\! 10^{-3}$ \\
Bern(alg)        & 180\,ms  & 1.05$\times$ & 0.00609 & $\checkmark$ & 24 & $1.5 \!\times\! 10^{-2}$ \\
\bottomrule
\end{tabular}
\end{table}

\textbf{주요 관측}:
\begin{itemize}
  \item 모든 방법이 수렴하며 비용 함수 $J$가 $0.006$--$0.007$ 범위로 일치한다.
  \item Affine(num)이 96\,ms로 최고속이며, RK4 대비 \textbf{1.97배} 빠르다.
  \item 대수적 보정도 RK4와 동등하거나 빠른 속도를 달성하였다
        (Affine: 1.33$\times$, Bern: 1.05$\times$).
  \item Bern(num)과 Bern(alg)의 BC 위반 차이
        ($3.3 \!\times\! 10^{-3}$ vs $1.5 \!\times\! 10^{-2}$)는
        보정 모드에 따른 참조 위치 정확도 차이를 반영한다.
\end{itemize}

\subsection{비용 구조 분석}

표~\ref{tab:cost_breakdown}은 SCP 1회 반복당 추정 비용 구성이다.

\begin{table}[h]
\centering
\caption{SCP 1회 반복당 비용 구성 (시나리오 B, 추정)}
\label{tab:cost_breakdown}
\begin{tabular}{lrrrr}
\toprule
구성요소 & RK4 & Affine(num) & Affine(alg) & Bern(alg) \\
\midrule
참조 궤적/중력 보정 & 4.0  & 2.4 & 2.5 & 3.5 \\
드리프트 적분       & 0.04 & 0.5 & 0.5 & 0.4 \\
QP/SOCP (CVXPY)     & 3.5  & 3.5 & 3.5 & 3.5 \\
\midrule
\textbf{반복당 합계 [ms]} & \textbf{7.5} & \textbf{6.4} & \textbf{6.5} & \textbf{7.4} \\
\bottomrule
\end{tabular}
\end{table}

워밍스타트 + 적응적 조기 종료 덕에 중력 보정 반복이
SCP 초기에는 3--4회, 수렴 근방에서는 1회로 줄어
전체 비용에서 보정이 차지하는 비중이 작아졌다.

\subsection{Numerical이 Algebraic보다 빠른 이유}

대수적 보정의 1회 반복은 \texttt{gravity\_composition\_pipeline}을 호출하여
\texttt{bern\_product}를 $\sim$20회 순차 실행한다.
각 호출은 벡터화되어 빠르지만, Python 함수 호출 + 배열 생성 오버헤드가
호출당 $\sim$0.02\,ms씩 누적되어 $\sim$0.4\,ms/반복이다.

반면 수치적 보정은 행렬곱 + 벡터 중력평가 + 누적합 + lstsq =
\textbf{5회}의 bulk numpy 호출로 $\sim$0.08\,ms/반복이다.

즉 연산량(FLOP)은 대수적이 적지만, \textbf{Python 레벨 함수 호출 횟수}에서
$\sim$4배 차이가 발생한다.


%% ══════════════════════════════════════════════════════════════
\section{최적화 구조적 함의}
\label{sec:structure}
%% ══════════════════════════════════════════════════════════════

연산 속도만으로는 수치적 보정이 우위이나,
제어점에 대한 볼록 최적화의 관점에서는 대수적 보정이
본질적으로 다른 가치를 제공한다.

\subsection{수치적 보정: $\mathbf{c}_v$는 불투명한 상수}

수치적 보정에서 드리프트 적분은 이산 점에서의 중력 평가와 수치적분을
거쳐 산출된다.
결과 $\mathbf{c}_v$, $\mathbf{c}_r$은 제어점 $\mathbf{P}_u$에 대한
함수적 관계가 소실된 ``불투명한 상수''로서 QP/SOCP에 투입된다.
따라서 매 SCP 반복마다 $\mathbf{c}_v$, $\mathbf{c}_r$을 새로 계산해야 하며,
이들의 제어점 감도 $\partial \mathbf{c}_v / \partial \mathbf{P}_u$를
구하려면 $3(N+1)$회의 유한차분이 필요하다.

\subsection{대수적 보정: $\mathbf{c}_v$는 $\mathbf{P}_u$의 명시적 함수}

대수적 보정에서는 전체 계산 체인이 제어점 공간의 대수적 연산으로 구성된다:
\[
  \mathbf{P}_u \;\xrightarrow{\text{선형}}\;
  \mathbf{q}_r \;\xrightarrow{\text{Bernstein 대수}}\;
  \mathbf{q}_{a} \;\xrightarrow{\text{선형 (mean)}}\;
  \mathbf{c}_v
\]

이 구조는 두 가지를 가능하게 한다.

\paragraph{(a) 커플링 행렬의 해석적 계산.}
보고서~009\cite{report009}의 커플링 행렬
\begin{equation}
  \mathbf{c}_v \approx \mathbf{c}_{v,\mathrm{ref}} + M_v \, \Delta\mathbf{P}_u
  \label{eq:coupling}
\end{equation}
을 해석적으로 산출할 수 있다.
이를 경계조건에 반영하면 $\mathbf{P}_u$에 대해 여전히 선형이면서도
더 정확한 선형화가 가능하여 SCP 수렴이 가속된다.

\paragraph{(b) 단일 볼록 정식화 가능성.}
모든 비선형성이 Bernstein 대수 (곱, 합성, 차수 축소)로 표현되면,
중력의 제어점 의존성을 경계조건에 명시적으로 포함하여
반복 선형화 없이 \textbf{단일 볼록 최적화}로 풀 수 있는 구조가 열린다.
이는 보고서~010\cite{report010}이 제시한 ``수치적분 없는 드리프트 적분''의
궁극적 활용 방향이다.

\subsection{요약: 속도 vs 구조}

\begin{table}[h]
\centering
\caption{수치적 vs 대수적 보정의 특성 비교}
\label{tab:num_vs_alg}
\begin{tabular}{lcc}
\toprule
특성 & Numerical & Algebraic \\
\midrule
SCP 1회 반복 비용        & 빠름 (5회 호출) & 느림 (20회 호출) \\
$\mathbf{c}_v$와 $\mathbf{P}_u$의 관계 & 불투명 (상수) & \textbf{명시적 함수} \\
커플링 행렬 $M_v$        & 유한차분 필요    & \textbf{해석적 계산 가능} \\
단일 볼록 정식화          & 불가능          & \textbf{가능성 있음} \\
\bottomrule
\end{tabular}
\end{table}

대수적 보정의 본질적 가치는 연산 속도가 아니라,
\textbf{제어점 최적화 구조를 보존한다}는 점에 있다.


%% ══════════════════════════════════════════════════════════════
\section{결론}
%% ══════════════════════════════════════════════════════════════

\begin{enumerate}
  \item 보고서 009, 010의 드리프트 적분 방법을 SCP 솔버에 통합하고,
        제어 전용 위치에 빠진 중력 기여를 자기일관 반복으로 보정하는
        방법을 제안·구현하였다.

  \item Python 루프의 벡터화, 가중치 행렬 캐싱, 워밍스타트, 적응적 조기 종료,
        보정용 $K$ 분리를 적용하여
        초기 구현 대비 \textbf{180배} 속도 향상을 달성하였다.

  \item 최종 벤치마크에서 Affine(numerical)은 RK4 대비 \textbf{1.97배},
        Affine(algebraic)은 \textbf{1.33배} 빠르면서
        모든 방법이 유사한 비용과 수렴을 보였다.

  \item 대수적 보정은 연산 속도에서는 수치적 보정에 뒤지지만,
        드리프트 적분의 제어점 의존성을 명시적으로 보존하여
        커플링 행렬 산출 및 단일 볼록 정식화로의 확장이 가능하다.
        이는 Bernstein 대수 프레임워크의 구조적 우위이다.
\end{enumerate}


%% ── 참고문헌 ─────────────────────────────────────────────────
\bibliographystyle{plain}
\bibliography{references}

\end{document}
